\section{宫门、遗诏与八王之乱}
\label{sec:western-jin-court-politics-eight-princes}

太熙末年，西晋所要处理的并不只是皇帝更替，而是谁能够把诏令、尚书省与宫城卫队接到同一条指挥线上。武帝病重时，\eventanchor{JH-0289-YANG-JUN-REGENT}{西晋·杨骏辅政}外戚杨骏受遗诏辅政；随后武帝去世，太子司马衷即位，是为惠帝，杨骏以太傅、录尚书事主持朝政。《晋书·武帝纪》《惠帝纪》与《杨骏传》共同保存了这一权力转换的大致骨架：皇位传给了司马氏，而中枢的日常裁决暂时集中到后族手中。遗诏究竟怎样措辞、武帝原本要让哪些大臣共参辅政、诏令又经何人传达，史料并没有同样可靠的答案，不能把后来胜者的说法补成一套无缝的临终安排。

杨骏的难处，在于他既要以顾命大臣的名义维持新朝廷，又无法把皇后贾氏、宗王和宫中武力排除在政治之外。惠帝本人能否独立裁决，传世叙事多有贬抑而难以直接检验；更可把握的是，任官、奏报和宫门的控制权已成为各方必须争取的资源。统一后的皇室封王本可分担地方防务，此时却也提供了能带兵入京、以宗室名义干预中枢的行动者。继承没有建立一套约束后族和宗王的程序，反而使顾命、后族与宫卫之间的临时均衡格外脆弱。

\subsection{杨氏倒台，宫廷政变成为先例}

永平元年，贾后联络楚王司马玮等发动政变，杨骏被杀，杨氏集团随之失势。这一结果见于《晋书·惠帝纪》《贾后传》《杨骏传》，《资治通鉴》则把它编入连续的朝政更替。两类叙述都说明，宗王和禁卫能够被动员进入宫廷冲突；但贾后、司马玮和各禁军将领各自如何谋划，宫中诏书是否尽为真传，不能由后出的叙事一一坐实。把政变简单归结为一人“阴谋”，会遮蔽更重要的条件：杨氏排斥异己，皇后缺少参与朝政的稳定位置，而持兵宗王已被引入解决中枢争端。

杨氏倒台没有恢复顾命政治。汝南王司马亮与卫瓘以辅政者身份出现，不久又在贾后借司马玮发动的清洗中被杀，司马玮旋即被处死。短时间内连续出现的诏命、捕杀和再处置，使“谁代皇帝发令”的问题本身成为斗争的一部分。《晋书》诸传对于先后诏令和责任归属并不完全一致，《资治通鉴》的整合也不等于获得了一份独立的宫中记录。可以确定的不是某一纸诏书的全部来历，而是辅政者无法凭官号自保，宗王带兵进入都城也从例外变成可被反复调用的办法。

\subsection{关陇的警报没有进入洛阳的优先次序}

宫门内反复改换主人时，关中已经显示出更难处理的压力。元康六年至九年，氐、羌集团中的齐万年起兵，西晋多次征讨，直到二九九年才大体平定。战事的延续、主要将领与齐万年败亡，可由《晋书·惠帝纪》以及《张方传》《江统传》互相参照；参与的部落究竟有哪些、兵员多少、关中民众分别支持何方，却不能从史书中的概称推得。史书所用的群体名称首先是当时官府和编史者的分类，不应直接变成整齐而固定的族群图像。

这场战争耗去关陇的军粮和机动兵力，也暴露朝廷缺乏能够长期处理边地安置、土地与军镇关系的办法。江统《徙戎论》留下了朝臣对关中局势的强烈忧惧；它是理解这种忧惧和排斥性政策语言的重要文本，却不是人口规模、迁徙方向或叛乱原因的统计报告。不能因此把齐万年之乱解释成单一族群的必然冲突。较稳的结论是，西部防务已经需要持久投入，而洛阳的权力竞争仍不断抽走能够统筹投入的官员、军队与注意力。边地的紧张不是宫廷斗争的背景布置，它正在削弱中枢重新建立秩序的余地。

\subsection{废太子与诸王相攻}

永康元年，贾后以谋反罪废杀太子司马遹，继承危机由宫内猜忌变成公开的政治裂缝。《晋书·愍怀太子遹传》《贾后传》记有伪造诏书、诱逼书写等情节，但这些细节的传递链和执行过程难以逐项复原。太子被废杀这一结果本身已足以改变局势：皇位继承失去被普遍承认的次序，反对贾后的宗王因而得到可以动员官僚和军队的名义。

次年，赵王司马伦杀贾后，废惠帝而自立为帝。以“清君侧”进入宫廷的行动，很快越过了恢复旧秩序的界线。司马伦的废立与称帝，可由《晋书·赵王伦传》和《惠帝纪》确证；诏册由谁起草、多少朝臣出于赞同而非胁迫、贾后死亡的具体情状，则不宜写成确定事实。宗王自称皇帝触犯了维系司马氏秩序的底线，齐王司马冏等随即起兵，复立惠帝，并以大司马辅政。惠帝回到皇位，并不表示皇权重新获得军权；它只说明一支宗王联盟暂时以复辟的名义取得了洛阳。

后人以“八王之乱”概括这一连串冲突，名称有助于辨认其规模，却不应误作参加者人数、阵营或战争阶段都完全固定的清单。司马冏掌政后，河间王司马颙与成都王司马颖起兵；三〇二年，长沙王司马乂在洛阳杀司马冏，接掌朝政。关中、邺城和洛阳之间的军队被接连调入中枢，使辅政从官僚事务变成依靠外镇兵力维持的职位。《晋书·齐王冏传》《长沙厉王乂传》《成都王颖传》所载各军协同和城内响应，不足以精确还原每一支军队的数额；但政权在数月间数度易手，已经说明中央不能自行限制其所倚赖的武装。

三〇三年，司马乂据洛阳抵抗司马颙、司马颖联军，京师陷入长期战斗。围城期间的粮价、死伤与惠帝是否亲自指挥，材料说法纷歧，不能把任何一个数字当作城市的完整经验。战斗持续本身却有明确的制度后果：皇帝、官署和城防被系于一方宗王的军事胜负，财政、漕运和地方治安则无从正常调度。巴蜀的李雄同年攻取成都，次年刘渊在左国城自称汉王，陈敏也在江东据地自置官属；这些并行的地方与边地变化不必被夸写为由洛阳一役直接造成，却显示中央内战发生时，原先由晋廷维系的区域节点已难得到有效支援。刘渊所能动员的部众范围、陈敏得到多少江东士族支持，分别见于《晋书·刘元海载记》《陈敏传》及《资治通鉴》《建康实录》的不同叙述，均不能据此画成边界明确、内部一致的政权。

永安元年，司马乂被东海王司马越等交给司马颙军并遭杀害，外军进入并支配洛阳。司马越何以倒戈、交付的程序为何、司马乂如何遇害，是《晋书·长沙厉王乂传》《东海孝献王越传》之间仍须保留的疑点；司马乂失去洛阳、惠帝被迫移置的结果则较为清楚。三〇五年，司马越又以讨伐司马颙为名起兵，更多州郡被卷入战争。各地响应的次序、兵额和地方社会被征发到何种程度，都不能由“勤王”或“讨逆”的称号推定；可以看见的是，原该守卫道路和边境的军队与粮道，持续服务于争夺皇帝所在地的内战。

\subsection{空出的皇位与空出的中枢}

三〇六年，惠帝去世，司马炽即位，是为怀帝，司马越仍掌握主要军政权。《晋书·惠帝纪》《怀帝纪》与《东海孝献王越传》可相互印证继位和权力格局，惠帝死因、司马越是否应负责任以及新帝能够实际决定多少政务，则不应强作结论。同年，李雄在成都称帝，说明西南的割据已形成自己的皇帝体制；关中、并州与江东的局势也都不再能仅靠洛阳的任命整合。怀帝所继承的是皇帝名号和残存官署，而不是一支可统一指挥、可稳定供给的国家力量。

因此，西晋的崩解不能只被解释为某位皇帝的个人才能，或某一位后妃、宗王的品行。杨骏辅政、贾后政变、废杀太子和诸王相攻，使继承和诏令先后失去可信的程序；齐万年之乱所显出的关陇压力、成都与左国城的政权变化，又使中枢无力把内战限定在宫城之内。到怀帝即位时，中央的破裂已与地方、边地和道路的失序相互牵连。此后洛阳所面对的危机，不能再靠一次复辟或一次辅政来化解。

\subsection{从内战到亡国}

怀帝继位后，空出的中枢很快面对一个新的皇帝中心。三〇八年，刘渊在平阳称帝，以汉为国号；其官制究竟在多大程度上仿汉、汉室继承语言如何被不同部众理解，不能由帝号本身断定，但称帝确实把并州、河东的多支武装放入了可以任官、征发的竞争框架。司马越仍握有主要军政权，却未能重建洛阳的外援。三一〇年他带兵东出途中去世，部众与辎重失去统一指挥，京师的粮道和可用机动力量更加空虚。

三一一年，汉赵军攻陷洛阳，怀帝被俘押往平阳。死散、杀掠和宫城破坏的数字，以及皇帝出宫和被俘后的礼仪，都有不同叙事，不能把它们合成一幅精确的灾难统计图；都城陷落、皇帝被俘和官僚中心破坏却足以说明，西晋已经失去维持全国任命与征发的支点。三一三年怀帝遇害，长安拥立司马邺为愍帝，是在残余官僚与宗室中保存晋朝继承线的努力，而不是对北方秩序的恢复。

三一六年，刘曜攻陷长安，愍帝出降。围城饥荒、投降条件与关中各地的反应不宜由一两种纪传写成定论；长安的失守和西晋北方朝廷的终结没有疑义。内战所耗散的军队、道路和政治信任，到此终于与汉赵的进攻连在一起。西晋的结局不是洛阳宫门中某一派战胜另一派的最后一幕，而是中枢已不能把皇统、粮运与地方军镇重新接合；晋室此后只能在建业另觅存续的条件。

\subsection{史书、编年与碑刻所能说明的不同问题}

本节的年序主要依靠《晋书》的帝纪和相关列传、载记。《晋书》由唐初官修，广采旧史与杂传，保存了西晋宫廷冲突、汉赵进攻与两京陷落最完整的传世叙事；也正因如此，其中的谶纬、人物褒贬、宫闱秘事和事后归责，必须同事件发生本身分开处理。《资治通鉴》把这些材料编入连续年序，便于辨认政变、征讨、继位与陷都之间的前后关系；它是北宋的编纂成果，有自己的删取和正统判断，不能被当作与《晋书》彼此独立的现场证词。两书相合，可以加强对废立、入京和战事结果的把握，却不足以证实诏书真伪、人物动机或兵员数字。

碑刻、墓志和地方考古提供的是另一种尺度：它们有时能为姓名、官职、纪日与地方网络提供独立锚点，却很少能够记录二八九至三一六年每一次宫廷密令、城门争夺或军队调动。本节没有以一方墓志或一通碑刻来证明杨骏、贾后、诸王或刘渊的谋划，也没有拿都城遗址推算全城人口与损失。它们提醒我们，官修史书以洛阳和宗室为中心的视野，不能代替关中军民、边地部众和被征发者的经历。正是在这些材料缺口中，叙事只能写明账本所列的继承、政变、战争与政权变化，并把无从核实的言辞、心理和场景留在史料之外。

% 深描扩写：资源、道路与人的选择
\section{在事件之间：资源、道路与人的选择}
\subsection{顾命与宫门}
从一件可见之物进入，才不会把这段历史写成抽象的王朝更替。 二八九至二九〇年杨骏受遗诏辅政，惠帝即位。洛阳宫城、尚书省与禁卫驻地中的遗诏、宫门、诏版和兵符不是背景，而是命令、粮食、人员和消息能够移动的条件。杨骏、惠帝、杨太后、贾后、顾命大臣被放在同一条不平坦的关系链上：有人借它发布号令，有人依它转运供给，也有人在它被截断时改变依附。顾命给外戚以名义，却不能排除皇后、宗王和掌握宫卫者对任命的介入。

若只把二八九至二九〇年杨骏受遗诏辅政，惠帝即位归结为一位皇帝、将领或谋臣的性格，便会错过更长的压力。继承安排、军镇分布、土地关系、边防开支和迁徙，往往早已改变了地方的承受能力；眼前的冲突只是把它们集中显露。能署名的诏令未必有军队，掌兵者也未必得到广泛承认。 因而名义、武力和供给必须并置叙述，不能把其中任何一项当成万能的解释。

从地方尺度看，事情并不只发生在帝纪最常书写的宫城或战场。守令、吏员、舟人、农户、工匠、商人和随军家属未必留下姓名，却要承担征发、迁移、停市、守城或重登簿册的后果。仓储是否有粮，渡口是否可用，旧籍能否辨认，新来者与旧居民能否取得暂时位置，决定命令能否成为日常秩序。

人事转折也应在这种约束中理解。杨骏、惠帝、杨太后、贾后、顾命大臣不是脱离制度而自由行动的人物，他们受官号、部众、家族、道路、既得任命以及对对手力量的有限消息所限。传记常把胜者写得果断、败者写得必然失败；它保存冲突轮廓，却容易遮住联盟为何在此刻成立，又为何在城防、补给或继承改变后瓦解。

二八九至二九〇年杨骏受遗诏辅政，惠帝即位的直接结果，通常不是把一片地方从混乱变为整齐。它会重写谁守备、哪些户口重登、哪条水道由谁巡护、哪些旧官仍可使用，也会让一部分人带着家属、技术、经卷或军籍离开原位。新获土地需要守军，新接军队需要粮饷，新立皇统需要文书；这些持续成本很快进入下一场危机，形成难以回避的反馈回路。

这也是制度得失必须落在具体人事上的原因。能够提高国家可见性的制度，可能让田地、户口与劳役更容易被比较和调动；它也可能把负担加在原本已受战乱、迁徙或边防牵累的家庭身上。地方中介既能使制度落地，也能借掌握信息、仓储或武力改变其样子。所谓成功，不宜只看诏令颁布或城池攻下，还要看这种接合能否持续。

本节年序主要依据《晋书》帝纪、列传与《资治通鉴》，统一时并参照《三国志》《建康实录》。这些材料足以确认二八九至二九〇年杨骏受遗诏辅政，惠帝即位及其前后关系，使顾命与宫门进入可核对的政治脊柱；但成书时间、编纂立场与保存条件并不相同。官修史书长于保存年序、政权语言与显贵行动，地方材料长于锚定特定人地；它们各有沉默，不能把诏令、兵数、族名或人物动机直接扩展成全社会的事实。 因此，叙述把能够互证的政权转换、行军、任命和城市易手写成事实，把兵数、伤亡、族属边界、地方响应和人物动机保留为需要继续检验的问题。

墓志、题记、文书、城址与道路遗存提供的是另一种尺度。它们有时能够校正姓名、日期、官职和特定的社会网络，也能让都城叙事之外的材料进入视野；它们却不能单独代表整座城市、一个州郡或所有流民。把这两类材料放在一起，不是为了制造一幅无缝全景，而是为了知道哪些判断可以成立，哪些地方仍应保留空白。

从这个角度回看，顾命与宫门不是供读者逐条勾选的事件卡，而是持续展开的历史处境。遗诏、宫门、诏版和兵符使国家能力可见，也使它的脆弱处可见：只要文书、军粮、人员和道路不能重新接合，皇帝名义便不能替代地方日常。后来的政权会以新制度、宗教安排、军事编制或城市工程继续回答这一难题，同时也会制造新的受益者、成本和不确定性。

\subsection{贾后政变的先例}
这里最重要的不是替后人判断成败，而是辨认当时有哪些资源已经断裂或仍可使用。 二九一年贾后联络宗王和禁军诛杀杨骏。洛阳城内的宫门和官署通道中的诏令、夜间城门和集结军队不是背景，而是命令、粮食、人员和消息能够移动的条件。贾后、司马玮、杨骏集团、禁卫将领被放在同一条不平坦的关系链上：有人借它发布号令，有人依它转运供给，也有人在它被截断时改变依附。宗王兵力进入宫廷并改变辅政者命运，改变了此后所有人对中枢的预期。

若只把二九一年贾后联络宗王和禁军诛杀杨骏归结为一位皇帝、将领或谋臣的性格，便会错过更长的压力。继承安排、军镇分布、土地关系、边防开支和迁徙，往往早已改变了地方的承受能力；眼前的冲突只是把它们集中显露。有人能守城，有人能运粮，也有人只剩下离开原处的道路。 因而名义、武力和供给必须并置叙述，不能把其中任何一项当成万能的解释。

从地方尺度看，事情并不只发生在帝纪最常书写的宫城或战场。守令、吏员、舟人、农户、工匠、商人和随军家属未必留下姓名，却要承担征发、迁移、停市、守城或重登簿册的后果。仓储是否有粮，渡口是否可用，旧籍能否辨认，新来者与旧居民能否取得暂时位置，决定命令能否成为日常秩序。

人事转折也应在这种约束中理解。贾后、司马玮、杨骏集团、禁卫将领不是脱离制度而自由行动的人物，他们受官号、部众、家族、道路、既得任命以及对对手力量的有限消息所限。传记常把胜者写得果断、败者写得必然失败；它保存冲突轮廓，却容易遮住联盟为何在此刻成立，又为何在城防、补给或继承改变后瓦解。

二九一年贾后联络宗王和禁军诛杀杨骏的直接结果，通常不是把一片地方从混乱变为整齐。它会重写谁守备、哪些户口重登、哪条水道由谁巡护、哪些旧官仍可使用，也会让一部分人带着家属、技术、经卷或军籍离开原位。新获土地需要守军，新接军队需要粮饷，新立皇统需要文书；这些持续成本很快进入下一场危机，形成难以回避的反馈回路。

这也是制度得失必须落在具体人事上的原因。能够提高国家可见性的制度，可能让田地、户口与劳役更容易被比较和调动；它也可能把负担加在原本已受战乱、迁徙或边防牵累的家庭身上。地方中介既能使制度落地，也能借掌握信息、仓储或武力改变其样子。所谓成功，不宜只看诏令颁布或城池攻下，还要看这种接合能否持续。

本节年序主要依据《晋书》帝纪、列传与《资治通鉴》，统一时并参照《三国志》《建康实录》。这些材料足以确认二九一年贾后联络宗王和禁军诛杀杨骏及其前后关系，使贾后政变的先例进入可核对的政治脊柱；但成书时间、编纂立场与保存条件并不相同。官修史书长于保存年序、政权语言与显贵行动，地方材料长于锚定特定人地；它们各有沉默，不能把诏令、兵数、族名或人物动机直接扩展成全社会的事实。 因此，叙述把能够互证的政权转换、行军、任命和城市易手写成事实，把兵数、伤亡、族属边界、地方响应和人物动机保留为需要继续检验的问题。

墓志、题记、文书、城址与道路遗存提供的是另一种尺度。它们有时能够校正姓名、日期、官职和特定的社会网络，也能让都城叙事之外的材料进入视野；它们却不能单独代表整座城市、一个州郡或所有流民。把这两类材料放在一起，不是为了制造一幅无缝全景，而是为了知道哪些判断可以成立，哪些地方仍应保留空白。

从这个角度回看，贾后政变的先例不是供读者逐条勾选的事件卡，而是持续展开的历史处境。诏令、夜间城门和集结军队使国家能力可见，也使它的脆弱处可见：只要文书、军粮、人员和道路不能重新接合，皇帝名义便不能替代地方日常。后来的政权会以新制度、宗教安排、军事编制或城市工程继续回答这一难题，同时也会制造新的受益者、成本和不确定性。

\subsection{齐万年与关中警报}
现存材料不能还原每一个现场，却能让人看见选择被空间和供给怎样限定。 二九六至二九九年齐万年起兵，西晋长期征讨。关中、雍州与西北道路中的军粮、山道、堡壁和征发文书不是背景，而是命令、粮食、人员和消息能够移动的条件。齐万年、氐羌部众、晋军、江统、当地编户被放在同一条不平坦的关系链上：有人借它发布号令，有人依它转运供给，也有人在它被截断时改变依附。边地安置、土地和军费已须持续投入，洛阳却仍消耗统筹能力。

若只把二九六至二九九年齐万年起兵，西晋长期征讨归结为一位皇帝、将领或谋臣的性格，便会错过更长的压力。继承安排、军镇分布、土地关系、边防开支和迁徙，往往早已改变了地方的承受能力；眼前的冲突只是把它们集中显露。政治语言可以称勤王、讨逆或归附，成本仍由持兵者、供给者和居民共同承担。 因而名义、武力和供给必须并置叙述，不能把其中任何一项当成万能的解释。

从地方尺度看，事情并不只发生在帝纪最常书写的宫城或战场。守令、吏员、舟人、农户、工匠、商人和随军家属未必留下姓名，却要承担征发、迁移、停市、守城或重登簿册的后果。仓储是否有粮，渡口是否可用，旧籍能否辨认，新来者与旧居民能否取得暂时位置，决定命令能否成为日常秩序。

人事转折也应在这种约束中理解。齐万年、氐羌部众、晋军、江统、当地编户不是脱离制度而自由行动的人物，他们受官号、部众、家族、道路、既得任命以及对对手力量的有限消息所限。传记常把胜者写得果断、败者写得必然失败；它保存冲突轮廓，却容易遮住联盟为何在此刻成立，又为何在城防、补给或继承改变后瓦解。

二九六至二九九年齐万年起兵，西晋长期征讨的直接结果，通常不是把一片地方从混乱变为整齐。它会重写谁守备、哪些户口重登、哪条水道由谁巡护、哪些旧官仍可使用，也会让一部分人带着家属、技术、经卷或军籍离开原位。新获土地需要守军，新接军队需要粮饷，新立皇统需要文书；这些持续成本很快进入下一场危机，形成难以回避的反馈回路。

这也是制度得失必须落在具体人事上的原因。能够提高国家可见性的制度，可能让田地、户口与劳役更容易被比较和调动；它也可能把负担加在原本已受战乱、迁徙或边防牵累的家庭身上。地方中介既能使制度落地，也能借掌握信息、仓储或武力改变其样子。所谓成功，不宜只看诏令颁布或城池攻下，还要看这种接合能否持续。

本节年序主要依据《晋书》帝纪、列传与《资治通鉴》，统一时并参照《三国志》《建康实录》。这些材料足以确认二九六至二九九年齐万年起兵，西晋长期征讨及其前后关系，使齐万年与关中警报进入可核对的政治脊柱；但成书时间、编纂立场与保存条件并不相同。官修史书长于保存年序、政权语言与显贵行动，地方材料长于锚定特定人地；它们各有沉默，不能把诏令、兵数、族名或人物动机直接扩展成全社会的事实。 因此，叙述把能够互证的政权转换、行军、任命和城市易手写成事实，把兵数、伤亡、族属边界、地方响应和人物动机保留为需要继续检验的问题。

墓志、题记、文书、城址与道路遗存提供的是另一种尺度。它们有时能够校正姓名、日期、官职和特定的社会网络，也能让都城叙事之外的材料进入视野；它们却不能单独代表整座城市、一个州郡或所有流民。把这两类材料放在一起，不是为了制造一幅无缝全景，而是为了知道哪些判断可以成立，哪些地方仍应保留空白。

从这个角度回看，齐万年与关中警报不是供读者逐条勾选的事件卡，而是持续展开的历史处境。军粮、山道、堡壁和征发文书使国家能力可见，也使它的脆弱处可见：只要文书、军粮、人员和道路不能重新接合，皇帝名义便不能替代地方日常。后来的政权会以新制度、宗教安排、军事编制或城市工程继续回答这一难题，同时也会制造新的受益者、成本和不确定性。

\subsection{太子被废}
把时间倒回当时，行动者面对的从来不是后来已经知道的结局。 三〇〇年贾后废杀太子司马遹。洛阳宫城及宗室官僚网络中的太子诏书、囚禁处所和传信者不是背景，而是命令、粮食、人员和消息能够移动的条件。贾后、太子、宗王、朝臣、宫廷执行者被放在同一条不平坦的关系链上：有人借它发布号令，有人依它转运供给，也有人在它被截断时改变依附。合法继承人被除去，使反对后族者取得动员军队和官僚的强大名义。

若只把三〇〇年贾后废杀太子司马遹归结为一位皇帝、将领或谋臣的性格，便会错过更长的压力。继承安排、军镇分布、土地关系、边防开支和迁徙，往往早已改变了地方的承受能力；眼前的冲突只是把它们集中显露。位置、资源、联盟和有限消息限定选择，未被记录的心理不应由叙述替他们补足。 因而名义、武力和供给必须并置叙述，不能把其中任何一项当成万能的解释。

从地方尺度看，事情并不只发生在帝纪最常书写的宫城或战场。守令、吏员、舟人、农户、工匠、商人和随军家属未必留下姓名，却要承担征发、迁移、停市、守城或重登簿册的后果。仓储是否有粮，渡口是否可用，旧籍能否辨认，新来者与旧居民能否取得暂时位置，决定命令能否成为日常秩序。

人事转折也应在这种约束中理解。贾后、太子、宗王、朝臣、宫廷执行者不是脱离制度而自由行动的人物，他们受官号、部众、家族、道路、既得任命以及对对手力量的有限消息所限。传记常把胜者写得果断、败者写得必然失败；它保存冲突轮廓，却容易遮住联盟为何在此刻成立，又为何在城防、补给或继承改变后瓦解。

三〇〇年贾后废杀太子司马遹的直接结果，通常不是把一片地方从混乱变为整齐。它会重写谁守备、哪些户口重登、哪条水道由谁巡护、哪些旧官仍可使用，也会让一部分人带着家属、技术、经卷或军籍离开原位。新获土地需要守军，新接军队需要粮饷，新立皇统需要文书；这些持续成本很快进入下一场危机，形成难以回避的反馈回路。

这也是制度得失必须落在具体人事上的原因。能够提高国家可见性的制度，可能让田地、户口与劳役更容易被比较和调动；它也可能把负担加在原本已受战乱、迁徙或边防牵累的家庭身上。地方中介既能使制度落地，也能借掌握信息、仓储或武力改变其样子。所谓成功，不宜只看诏令颁布或城池攻下，还要看这种接合能否持续。

本节年序主要依据《晋书》帝纪、列传与《资治通鉴》，统一时并参照《三国志》《建康实录》。这些材料足以确认三〇〇年贾后废杀太子司马遹及其前后关系，使太子被废进入可核对的政治脊柱；但成书时间、编纂立场与保存条件并不相同。官修史书长于保存年序、政权语言与显贵行动，地方材料长于锚定特定人地；它们各有沉默，不能把诏令、兵数、族名或人物动机直接扩展成全社会的事实。 因此，叙述把能够互证的政权转换、行军、任命和城市易手写成事实，把兵数、伤亡、族属边界、地方响应和人物动机保留为需要继续检验的问题。

墓志、题记、文书、城址与道路遗存提供的是另一种尺度。它们有时能够校正姓名、日期、官职和特定的社会网络，也能让都城叙事之外的材料进入视野；它们却不能单独代表整座城市、一个州郡或所有流民。把这两类材料放在一起，不是为了制造一幅无缝全景，而是为了知道哪些判断可以成立，哪些地方仍应保留空白。

从这个角度回看，太子被废不是供读者逐条勾选的事件卡，而是持续展开的历史处境。太子诏书、囚禁处所和传信者使国家能力可见，也使它的脆弱处可见：只要文书、军粮、人员和道路不能重新接合，皇帝名义便不能替代地方日常。后来的政权会以新制度、宗教安排、军事编制或城市工程继续回答这一难题，同时也会制造新的受益者、成本和不确定性。

\subsection{八王争洛阳}
从一件可见之物进入，才不会把这段历史写成抽象的王朝更替。 三〇一至三〇五年诸王反复起兵、围攻与废立。洛阳、邺城、关中与入京道路中的城门、粮仓、诏册和随军粮车不是背景，而是命令、粮食、人员和消息能够移动的条件。司马伦、司马冏、司马颖、司马颙、司马越、司马乂被放在同一条不平坦的关系链上：有人借它发布号令，有人依它转运供给，也有人在它被截断时改变依附。皇帝所在地成为耗尽各地机动力量的磁石，外镇军力反而决定中枢。

若只把三〇一至三〇五年诸王反复起兵、围攻与废立归结为一位皇帝、将领或谋臣的性格，便会错过更长的压力。继承安排、军镇分布、土地关系、边防开支和迁徙，往往早已改变了地方的承受能力；眼前的冲突只是把它们集中显露。能署名的诏令未必有军队，掌兵者也未必得到广泛承认。 因而名义、武力和供给必须并置叙述，不能把其中任何一项当成万能的解释。

从地方尺度看，事情并不只发生在帝纪最常书写的宫城或战场。守令、吏员、舟人、农户、工匠、商人和随军家属未必留下姓名，却要承担征发、迁移、停市、守城或重登簿册的后果。仓储是否有粮，渡口是否可用，旧籍能否辨认，新来者与旧居民能否取得暂时位置，决定命令能否成为日常秩序。

人事转折也应在这种约束中理解。司马伦、司马冏、司马颖、司马颙、司马越、司马乂不是脱离制度而自由行动的人物，他们受官号、部众、家族、道路、既得任命以及对对手力量的有限消息所限。传记常把胜者写得果断、败者写得必然失败；它保存冲突轮廓，却容易遮住联盟为何在此刻成立，又为何在城防、补给或继承改变后瓦解。

三〇一至三〇五年诸王反复起兵、围攻与废立的直接结果，通常不是把一片地方从混乱变为整齐。它会重写谁守备、哪些户口重登、哪条水道由谁巡护、哪些旧官仍可使用，也会让一部分人带着家属、技术、经卷或军籍离开原位。新获土地需要守军，新接军队需要粮饷，新立皇统需要文书；这些持续成本很快进入下一场危机，形成难以回避的反馈回路。

这也是制度得失必须落在具体人事上的原因。能够提高国家可见性的制度，可能让田地、户口与劳役更容易被比较和调动；它也可能把负担加在原本已受战乱、迁徙或边防牵累的家庭身上。地方中介既能使制度落地，也能借掌握信息、仓储或武力改变其样子。所谓成功，不宜只看诏令颁布或城池攻下，还要看这种接合能否持续。

本节年序主要依据《晋书》帝纪、列传与《资治通鉴》，统一时并参照《三国志》《建康实录》。这些材料足以确认三〇一至三〇五年诸王反复起兵、围攻与废立及其前后关系，使八王争洛阳进入可核对的政治脊柱；但成书时间、编纂立场与保存条件并不相同。官修史书长于保存年序、政权语言与显贵行动，地方材料长于锚定特定人地；它们各有沉默，不能把诏令、兵数、族名或人物动机直接扩展成全社会的事实。 因此，叙述把能够互证的政权转换、行军、任命和城市易手写成事实，把兵数、伤亡、族属边界、地方响应和人物动机保留为需要继续检验的问题。

墓志、题记、文书、城址与道路遗存提供的是另一种尺度。它们有时能够校正姓名、日期、官职和特定的社会网络，也能让都城叙事之外的材料进入视野；它们却不能单独代表整座城市、一个州郡或所有流民。把这两类材料放在一起，不是为了制造一幅无缝全景，而是为了知道哪些判断可以成立，哪些地方仍应保留空白。

从这个角度回看，八王争洛阳不是供读者逐条勾选的事件卡，而是持续展开的历史处境。城门、粮仓、诏册和随军粮车使国家能力可见，也使它的脆弱处可见：只要文书、军粮、人员和道路不能重新接合，皇帝名义便不能替代地方日常。后来的政权会以新制度、宗教安排、军事编制或城市工程继续回答这一难题，同时也会制造新的受益者、成本和不确定性。

\subsection{刘渊与两京失守}
这里最重要的不是替后人判断成败，而是辨认当时有哪些资源已经断裂或仍可使用。 三〇四至三一六年刘渊建汉，洛阳、长安先后陷落。左国城、平阳、洛阳、长安与河东关中道路中的牧地、城门、逃亡队伍和断裂粮道不是背景，而是命令、粮食、人员和消息能够移动的条件。刘渊、刘曜、怀帝、愍帝、残余晋臣、城中居民被放在同一条不平坦的关系链上：有人借它发布号令，有人依它转运供给，也有人在它被截断时改变依附。中枢不能再把皇统、军粮与地方守备接合，晋室只得另觅存续条件。

若只把三〇四至三一六年刘渊建汉，洛阳、长安先后陷落归结为一位皇帝、将领或谋臣的性格，便会错过更长的压力。继承安排、军镇分布、土地关系、边防开支和迁徙，往往早已改变了地方的承受能力；眼前的冲突只是把它们集中显露。有人能守城，有人能运粮，也有人只剩下离开原处的道路。 因而名义、武力和供给必须并置叙述，不能把其中任何一项当成万能的解释。

从地方尺度看，事情并不只发生在帝纪最常书写的宫城或战场。守令、吏员、舟人、农户、工匠、商人和随军家属未必留下姓名，却要承担征发、迁移、停市、守城或重登簿册的后果。仓储是否有粮，渡口是否可用，旧籍能否辨认，新来者与旧居民能否取得暂时位置，决定命令能否成为日常秩序。

人事转折也应在这种约束中理解。刘渊、刘曜、怀帝、愍帝、残余晋臣、城中居民不是脱离制度而自由行动的人物，他们受官号、部众、家族、道路、既得任命以及对对手力量的有限消息所限。传记常把胜者写得果断、败者写得必然失败；它保存冲突轮廓，却容易遮住联盟为何在此刻成立，又为何在城防、补给或继承改变后瓦解。

三〇四至三一六年刘渊建汉，洛阳、长安先后陷落的直接结果，通常不是把一片地方从混乱变为整齐。它会重写谁守备、哪些户口重登、哪条水道由谁巡护、哪些旧官仍可使用，也会让一部分人带着家属、技术、经卷或军籍离开原位。新获土地需要守军，新接军队需要粮饷，新立皇统需要文书；这些持续成本很快进入下一场危机，形成难以回避的反馈回路。

这也是制度得失必须落在具体人事上的原因。能够提高国家可见性的制度，可能让田地、户口与劳役更容易被比较和调动；它也可能把负担加在原本已受战乱、迁徙或边防牵累的家庭身上。地方中介既能使制度落地，也能借掌握信息、仓储或武力改变其样子。所谓成功，不宜只看诏令颁布或城池攻下，还要看这种接合能否持续。

本节年序主要依据《晋书》帝纪、列传与《资治通鉴》，统一时并参照《三国志》《建康实录》。这些材料足以确认三〇四至三一六年刘渊建汉，洛阳、长安先后陷落及其前后关系，使刘渊与两京失守进入可核对的政治脊柱；但成书时间、编纂立场与保存条件并不相同。官修史书长于保存年序、政权语言与显贵行动，地方材料长于锚定特定人地；它们各有沉默，不能把诏令、兵数、族名或人物动机直接扩展成全社会的事实。 因此，叙述把能够互证的政权转换、行军、任命和城市易手写成事实，把兵数、伤亡、族属边界、地方响应和人物动机保留为需要继续检验的问题。

墓志、题记、文书、城址与道路遗存提供的是另一种尺度。它们有时能够校正姓名、日期、官职和特定的社会网络，也能让都城叙事之外的材料进入视野；它们却不能单独代表整座城市、一个州郡或所有流民。把这两类材料放在一起，不是为了制造一幅无缝全景，而是为了知道哪些判断可以成立，哪些地方仍应保留空白。

从这个角度回看，刘渊与两京失守不是供读者逐条勾选的事件卡，而是持续展开的历史处境。牧地、城门、逃亡队伍和断裂粮道使国家能力可见，也使它的脆弱处可见：只要文书、军粮、人员和道路不能重新接合，皇帝名义便不能替代地方日常。后来的政权会以新制度、宗教安排、军事编制或城市工程继续回答这一难题，同时也会制造新的受益者、成本和不确定性。

\begin{causalchain}[本节制度得失]
这些事件显示，最高政权的转换只能提供重新组织的起点。能否把官员、户籍、军队、仓储、宗教机构与地方社会重新接合，才决定新秩序能维持多久；这种接合的成本和不均衡，也会成为下一次危机的条件。
\end{causalchain}

% 深描续写：制度、社会与材料的连续尺度
\section{继承程序怎样让外镇进入都城}
第一，时间上的连续并不等于政治上的连续。一个年号被替换，或一座城的城门被打开，并不自动让原有的债务、亲族、服役关系和地方记忆消失；它们常在新政权的名目下继续运作，也会在新的征发中显出抵触。 在洛阳、关中、邺城、平阳和入京道路这一组相互连通又彼此牵制的地点，遗诏、兵符、城门、粮仓、马匹与传令文书把抽象的权力落到可以触摸的层面。杨骏辅政、贾后政变、太子被废、八王争战与两京失守并非孤立的名词串，而是让后族、宗王、禁军、外镇将领、边地部众、城中官民不得不在资源有限、消息不全的条件下重新排列位置的过程。

就继承程序怎样让外镇进入都城而言，关键不在于把每个变化归还给一个唯一原因，而在于追问它怎样与前后事件相接。洛阳失守后，晋室必须依赖南方重新保存皇统 这种前后相承不是机械决定论：同样的制度或道路，在不同的都城、边境和乡里可能产生不同结果，因为掌握中介、承担成本和接收消息的人并不相同。

年序与主要行动可由《晋书》帝纪列传、《三国志》《建康实录》与《资治通鉴》相互校读；但史书中的兵数、死亡、动机、族名和地方响应经常带有修辞或编纂立场。墓志、题记、文书、城址和交通遗存可以使具体姓名、日期与物质网络更清楚，却同样不能代替社会全景。因而这一节只在材料足以支撑之处推进解释，在无法互证之处保留疑问。

第二，空间必须被当作历史行动者。河流、关隘、山道、港口和城市并不只规定军队能否通过，也规定消息怎样迟到、粮食怎样损耗、逃亡者在哪里停留，以及中央命令在何处必须转交给地方中介。 在洛阳、关中、邺城、平阳和入京道路这一组相互连通又彼此牵制的地点，遗诏、兵符、城门、粮仓、马匹与传令文书把抽象的权力落到可以触摸的层面。杨骏辅政、贾后政变、太子被废、八王争战与两京失守并非孤立的名词串，而是让后族、宗王、禁军、外镇将领、边地部众、城中官民不得不在资源有限、消息不全的条件下重新排列位置的过程。

就继承程序怎样让外镇进入都城而言，关键不在于把每个变化归还给一个唯一原因，而在于追问它怎样与前后事件相接。洛阳失守后，晋室必须依赖南方重新保存皇统 这种前后相承不是机械决定论：同样的制度或道路，在不同的都城、边境和乡里可能产生不同结果，因为掌握中介、承担成本和接收消息的人并不相同。

年序与主要行动可由《晋书》帝纪列传、《三国志》《建康实录》与《资治通鉴》相互校读；但史书中的兵数、死亡、动机、族名和地方响应经常带有修辞或编纂立场。墓志、题记、文书、城址和交通遗存可以使具体姓名、日期与物质网络更清楚，却同样不能代替社会全景。因而这一节只在材料足以支撑之处推进解释，在无法互证之处保留疑问。

第三，行政的力量来自可重复的登记，而非一次壮观的宣示。官府要把人口、土地、军役和税粮写入可以核验的簿册，便须依靠识字的吏员、熟悉乡里的中介和仍能维持的仓储道路；这些条件缺一，制度语言就会在地方变形。 在洛阳、关中、邺城、平阳和入京道路这一组相互连通又彼此牵制的地点，遗诏、兵符、城门、粮仓、马匹与传令文书把抽象的权力落到可以触摸的层面。杨骏辅政、贾后政变、太子被废、八王争战与两京失守并非孤立的名词串，而是让后族、宗王、禁军、外镇将领、边地部众、城中官民不得不在资源有限、消息不全的条件下重新排列位置的过程。

就继承程序怎样让外镇进入都城而言，关键不在于把每个变化归还给一个唯一原因，而在于追问它怎样与前后事件相接。洛阳失守后，晋室必须依赖南方重新保存皇统 这种前后相承不是机械决定论：同样的制度或道路，在不同的都城、边境和乡里可能产生不同结果，因为掌握中介、承担成本和接收消息的人并不相同。

年序与主要行动可由《晋书》帝纪列传、《三国志》《建康实录》与《资治通鉴》相互校读；但史书中的兵数、死亡、动机、族名和地方响应经常带有修辞或编纂立场。墓志、题记、文书、城址和交通遗存可以使具体姓名、日期与物质网络更清楚，却同样不能代替社会全景。因而这一节只在材料足以支撑之处推进解释，在无法互证之处保留疑问。

第四，战争中的普通人并非只作为数字出现。即使材料很少记录他们的名字，征发、避难、停市、重新登记、失去耕作时间和寻找新的保护者，仍然是政治变动得以发生而又付出代价的方式。叙事不能替他们虚构心声，却应当保留这种成本。 在洛阳、关中、邺城、平阳和入京道路这一组相互连通又彼此牵制的地点，遗诏、兵符、城门、粮仓、马匹与传令文书把抽象的权力落到可以触摸的层面。杨骏辅政、贾后政变、太子被废、八王争战与两京失守并非孤立的名词串，而是让后族、宗王、禁军、外镇将领、边地部众、城中官民不得不在资源有限、消息不全的条件下重新排列位置的过程。

就继承程序怎样让外镇进入都城而言，关键不在于把每个变化归还给一个唯一原因，而在于追问它怎样与前后事件相接。洛阳失守后，晋室必须依赖南方重新保存皇统 这种前后相承不是机械决定论：同样的制度或道路，在不同的都城、边境和乡里可能产生不同结果，因为掌握中介、承担成本和接收消息的人并不相同。

年序与主要行动可由《晋书》帝纪列传、《三国志》《建康实录》与《资治通鉴》相互校读；但史书中的兵数、死亡、动机、族名和地方响应经常带有修辞或编纂立场。墓志、题记、文书、城址和交通遗存可以使具体姓名、日期与物质网络更清楚，却同样不能代替社会全景。因而这一节只在材料足以支撑之处推进解释，在无法互证之处保留疑问。

第五，精英集团也不是可以用单一标签概括的整体。宗室、后族、侨姓、旧族、军将、僧团和地方豪强都在不同地点拥有不一样的资源；他们可以在一场危机中结盟，下一场危机又因任命、土地或供给而分裂。 在洛阳、关中、邺城、平阳和入京道路这一组相互连通又彼此牵制的地点，遗诏、兵符、城门、粮仓、马匹与传令文书把抽象的权力落到可以触摸的层面。杨骏辅政、贾后政变、太子被废、八王争战与两京失守并非孤立的名词串，而是让后族、宗王、禁军、外镇将领、边地部众、城中官民不得不在资源有限、消息不全的条件下重新排列位置的过程。

就继承程序怎样让外镇进入都城而言，关键不在于把每个变化归还给一个唯一原因，而在于追问它怎样与前后事件相接。洛阳失守后，晋室必须依赖南方重新保存皇统 这种前后相承不是机械决定论：同样的制度或道路，在不同的都城、边境和乡里可能产生不同结果，因为掌握中介、承担成本和接收消息的人并不相同。

年序与主要行动可由《晋书》帝纪列传、《三国志》《建康实录》与《资治通鉴》相互校读；但史书中的兵数、死亡、动机、族名和地方响应经常带有修辞或编纂立场。墓志、题记、文书、城址和交通遗存可以使具体姓名、日期与物质网络更清楚，却同样不能代替社会全景。因而这一节只在材料足以支撑之处推进解释，在无法互证之处保留疑问。

第六，军事胜利的后果总要经过安置。俘虏、降军、旧吏、流民和被征服地区的仓储，并不会因为战报写下胜利就自动变为新王朝的力量；如何给粮、编伍、授官、迁置或准其留居，往往决定征服能否持续。 在洛阳、关中、邺城、平阳和入京道路这一组相互连通又彼此牵制的地点，遗诏、兵符、城门、粮仓、马匹与传令文书把抽象的权力落到可以触摸的层面。杨骏辅政、贾后政变、太子被废、八王争战与两京失守并非孤立的名词串，而是让后族、宗王、禁军、外镇将领、边地部众、城中官民不得不在资源有限、消息不全的条件下重新排列位置的过程。

就继承程序怎样让外镇进入都城而言，关键不在于把每个变化归还给一个唯一原因，而在于追问它怎样与前后事件相接。洛阳失守后，晋室必须依赖南方重新保存皇统 这种前后相承不是机械决定论：同样的制度或道路，在不同的都城、边境和乡里可能产生不同结果，因为掌握中介、承担成本和接收消息的人并不相同。

年序与主要行动可由《晋书》帝纪列传、《三国志》《建康实录》与《资治通鉴》相互校读；但史书中的兵数、死亡、动机、族名和地方响应经常带有修辞或编纂立场。墓志、题记、文书、城址和交通遗存可以使具体姓名、日期与物质网络更清楚，却同样不能代替社会全景。因而这一节只在材料足以支撑之处推进解释，在无法互证之处保留疑问。

第七，财政不只是一项抽象收入。谷帛、田租、劳役、船只、马匹、器械和工匠的移动，决定一个都城或军镇能把多少命令变成现实；当这些资源被宫廷争夺或远征吸走时，边地和乡里会先感到秩序变薄。 在洛阳、关中、邺城、平阳和入京道路这一组相互连通又彼此牵制的地点，遗诏、兵符、城门、粮仓、马匹与传令文书把抽象的权力落到可以触摸的层面。杨骏辅政、贾后政变、太子被废、八王争战与两京失守并非孤立的名词串，而是让后族、宗王、禁军、外镇将领、边地部众、城中官民不得不在资源有限、消息不全的条件下重新排列位置的过程。

就继承程序怎样让外镇进入都城而言，关键不在于把每个变化归还给一个唯一原因，而在于追问它怎样与前后事件相接。洛阳失守后，晋室必须依赖南方重新保存皇统 这种前后相承不是机械决定论：同样的制度或道路，在不同的都城、边境和乡里可能产生不同结果，因为掌握中介、承担成本和接收消息的人并不相同。

年序与主要行动可由《晋书》帝纪列传、《三国志》《建康实录》与《资治通鉴》相互校读；但史书中的兵数、死亡、动机、族名和地方响应经常带有修辞或编纂立场。墓志、题记、文书、城址和交通遗存可以使具体姓名、日期与物质网络更清楚，却同样不能代替社会全景。因而这一节只在材料足以支撑之处推进解释，在无法互证之处保留疑问。

第八，宗教材料让国家之外的组织变得可见。寺院、道观、译场、石窟和造像把施主、工匠、经卷、田产和交通连在一起；王室的赞助或禁断当然重要，但不能替代对地方信众、财物和营造节奏的具体追问。 在洛阳、关中、邺城、平阳和入京道路这一组相互连通又彼此牵制的地点，遗诏、兵符、城门、粮仓、马匹与传令文书把抽象的权力落到可以触摸的层面。杨骏辅政、贾后政变、太子被废、八王争战与两京失守并非孤立的名词串，而是让后族、宗王、禁军、外镇将领、边地部众、城中官民不得不在资源有限、消息不全的条件下重新排列位置的过程。

就继承程序怎样让外镇进入都城而言，关键不在于把每个变化归还给一个唯一原因，而在于追问它怎样与前后事件相接。洛阳失守后，晋室必须依赖南方重新保存皇统 这种前后相承不是机械决定论：同样的制度或道路，在不同的都城、边境和乡里可能产生不同结果，因为掌握中介、承担成本和接收消息的人并不相同。

年序与主要行动可由《晋书》帝纪列传、《三国志》《建康实录》与《资治通鉴》相互校读；但史书中的兵数、死亡、动机、族名和地方响应经常带有修辞或编纂立场。墓志、题记、文书、城址和交通遗存可以使具体姓名、日期与物质网络更清楚，却同样不能代替社会全景。因而这一节只在材料足以支撑之处推进解释，在无法互证之处保留疑问。

第九，族称与地域称谓首先是当时军政和书写的分类。它们可以帮助辨认征发、赏罚、迁徙或外交发生的语境，却不应被直接转换成内部同质、永恒不变的现代群体；不同材料中的同一个词，也可能具有不同边界。 在洛阳、关中、邺城、平阳和入京道路这一组相互连通又彼此牵制的地点，遗诏、兵符、城门、粮仓、马匹与传令文书把抽象的权力落到可以触摸的层面。杨骏辅政、贾后政变、太子被废、八王争战与两京失守并非孤立的名词串，而是让后族、宗王、禁军、外镇将领、边地部众、城中官民不得不在资源有限、消息不全的条件下重新排列位置的过程。

就继承程序怎样让外镇进入都城而言，关键不在于把每个变化归还给一个唯一原因，而在于追问它怎样与前后事件相接。洛阳失守后，晋室必须依赖南方重新保存皇统 这种前后相承不是机械决定论：同样的制度或道路，在不同的都城、边境和乡里可能产生不同结果，因为掌握中介、承担成本和接收消息的人并不相同。

年序与主要行动可由《晋书》帝纪列传、《三国志》《建康实录》与《资治通鉴》相互校读；但史书中的兵数、死亡、动机、族名和地方响应经常带有修辞或编纂立场。墓志、题记、文书、城址和交通遗存可以使具体姓名、日期与物质网络更清楚，却同样不能代替社会全景。因而这一节只在材料足以支撑之处推进解释，在无法互证之处保留疑问。

第十，因果关系必须保留层次。结构性压力可能长期存在，近期触发也许只是一次废立、叛乱或失守；使行动成为可能的往往是道路、仓储、联盟和信息。把这些层次分开，才能避免把时间上相邻的变化误写成唯一原因。 在洛阳、关中、邺城、平阳和入京道路这一组相互连通又彼此牵制的地点，遗诏、兵符、城门、粮仓、马匹与传令文书把抽象的权力落到可以触摸的层面。杨骏辅政、贾后政变、太子被废、八王争战与两京失守并非孤立的名词串，而是让后族、宗王、禁军、外镇将领、边地部众、城中官民不得不在资源有限、消息不全的条件下重新排列位置的过程。

就继承程序怎样让外镇进入都城而言，关键不在于把每个变化归还给一个唯一原因，而在于追问它怎样与前后事件相接。洛阳失守后，晋室必须依赖南方重新保存皇统 这种前后相承不是机械决定论：同样的制度或道路，在不同的都城、边境和乡里可能产生不同结果，因为掌握中介、承担成本和接收消息的人并不相同。

年序与主要行动可由《晋书》帝纪列传、《三国志》《建康实录》与《资治通鉴》相互校读；但史书中的兵数、死亡、动机、族名和地方响应经常带有修辞或编纂立场。墓志、题记、文书、城址和交通遗存可以使具体姓名、日期与物质网络更清楚，却同样不能代替社会全景。因而这一节只在材料足以支撑之处推进解释，在无法互证之处保留疑问。

第十一，人物的选择值得叙述，却不应被神化。史料能让我们看见他们起兵、入城、纳降、逃奔或颁令，也能观察行动造成的后果；对未记录的密谈和心思，最诚实的写法是说明证据到此为止。 在洛阳、关中、邺城、平阳和入京道路这一组相互连通又彼此牵制的地点，遗诏、兵符、城门、粮仓、马匹与传令文书把抽象的权力落到可以触摸的层面。杨骏辅政、贾后政变、太子被废、八王争战与两京失守并非孤立的名词串，而是让后族、宗王、禁军、外镇将领、边地部众、城中官民不得不在资源有限、消息不全的条件下重新排列位置的过程。

就继承程序怎样让外镇进入都城而言，关键不在于把每个变化归还给一个唯一原因，而在于追问它怎样与前后事件相接。洛阳失守后，晋室必须依赖南方重新保存皇统 这种前后相承不是机械决定论：同样的制度或道路，在不同的都城、边境和乡里可能产生不同结果，因为掌握中介、承担成本和接收消息的人并不相同。

年序与主要行动可由《晋书》帝纪列传、《三国志》《建康实录》与《资治通鉴》相互校读；但史书中的兵数、死亡、动机、族名和地方响应经常带有修辞或编纂立场。墓志、题记、文书、城址和交通遗存可以使具体姓名、日期与物质网络更清楚，却同样不能代替社会全景。因而这一节只在材料足以支撑之处推进解释，在无法互证之处保留疑问。

第十二，后来的史书总在已知结局的立场上整理早期事件。它们提供不可替代的年序和政治词汇，也容易把成功者的选择写成必然。用同时代或地方材料提出反问，不是取消正史，而是避免让单一视角垄断问题。 在洛阳、关中、邺城、平阳和入京道路这一组相互连通又彼此牵制的地点，遗诏、兵符、城门、粮仓、马匹与传令文书把抽象的权力落到可以触摸的层面。杨骏辅政、贾后政变、太子被废、八王争战与两京失守并非孤立的名词串，而是让后族、宗王、禁军、外镇将领、边地部众、城中官民不得不在资源有限、消息不全的条件下重新排列位置的过程。

就继承程序怎样让外镇进入都城而言，关键不在于把每个变化归还给一个唯一原因，而在于追问它怎样与前后事件相接。洛阳失守后，晋室必须依赖南方重新保存皇统 这种前后相承不是机械决定论：同样的制度或道路，在不同的都城、边境和乡里可能产生不同结果，因为掌握中介、承担成本和接收消息的人并不相同。

年序与主要行动可由《晋书》帝纪列传、《三国志》《建康实录》与《资治通鉴》相互校读；但史书中的兵数、死亡、动机、族名和地方响应经常带有修辞或编纂立场。墓志、题记、文书、城址和交通遗存可以使具体姓名、日期与物质网络更清楚，却同样不能代替社会全景。因而这一节只在材料足以支撑之处推进解释，在无法互证之处保留疑问。

第十三，制度得失不能只用道德词评价。一个办法可能提高征发、守备或城市管理的能力，同时把负担转给某些家庭、军户或地方中介；这种副作用并不证明制度毫无作用，却解释了它为什么会在新的环境中引发反弹。 在洛阳、关中、邺城、平阳和入京道路这一组相互连通又彼此牵制的地点，遗诏、兵符、城门、粮仓、马匹与传令文书把抽象的权力落到可以触摸的层面。杨骏辅政、贾后政变、太子被废、八王争战与两京失守并非孤立的名词串，而是让后族、宗王、禁军、外镇将领、边地部众、城中官民不得不在资源有限、消息不全的条件下重新排列位置的过程。

就继承程序怎样让外镇进入都城而言，关键不在于把每个变化归还给一个唯一原因，而在于追问它怎样与前后事件相接。洛阳失守后，晋室必须依赖南方重新保存皇统 这种前后相承不是机械决定论：同样的制度或道路，在不同的都城、边境和乡里可能产生不同结果，因为掌握中介、承担成本和接收消息的人并不相同。

年序与主要行动可由《晋书》帝纪列传、《三国志》《建康实录》与《资治通鉴》相互校读；但史书中的兵数、死亡、动机、族名和地方响应经常带有修辞或编纂立场。墓志、题记、文书、城址和交通遗存可以使具体姓名、日期与物质网络更清楚，却同样不能代替社会全景。因而这一节只在材料足以支撑之处推进解释，在无法互证之处保留疑问。

\begin{sourcenote}[阅读这一连续尺度]
本节不把制度、宗教、战争和迁徙拆成彼此无关的栏目。继承程序怎样让外镇进入都城所涉及的事件，均以正史与编年材料维持最低年序，再以地方性材料限定结论范围。来源的差异不是缺陷，而是提醒读者：国家如何书写自身、地方留下何种痕迹、后来史家怎样重排叙事，回答的是不同问题。
\end{sourcenote}
